\documentclass[addpoints,12pt,answers]{exam}
\usepackage{babel}
\usepackage[top=3cm,bottom=3cm,left=2.5cm,right=2.5cm]{geometry}
\usepackage[utf8]{inputenc}
\usepackage[T1]{fontenc}
\usepackage{booktabs}
\usepackage{graphicx}
\usepackage{csquotes}
\usepackage{paralist}
%\usepackage{wasysym}
\usepackage[math]{iwona}
\usepackage{setspace}
\usepackage{enumitem}
\usepackage{caption}
\usepackage{subcaption}
\usepackage{amsmath}

\pointpoints{Point}{Points}
\bonuspointpoints{Bonuspunkt}{Bonuspunkte}
\renewcommand{\solutiontitle}{\noindent\textbf{Solution:}\enspace}

\chqword{Frage}   
\chpgword{Seite} 
\chpword{Punkte}   
\chbpword{Bonus Points} 
\chsword{Erreicht}   
\chtword{Gesamt}

\checkboxchar{\Square}
\checkedchar{\CheckedBox}

\pagestyle{headandfoot}
%\runningheadrule
%\firstpageheader{ECO208}{Test 3}{28th Feb, 2022}
%\runningheader{ECO208}{Test 3}{28th Feb, 2022}
%\firstpagefooter{ECO208}{Test 3}{\thepage\,/\,\numpages}
%\runningfooter{ECO208}{Test 3}{\thepage\,/\,\numpages}
\doublespacing

\begin{document}
	\vspace*{3em}
	
	\makebox[\textwidth]{Extra Question W3b}
	\makebox[\textwidth]{Name:\enspace\hrulefill}
	
	\vspace*{3em}
	
	\begin{questions}
		
		\question You are given the following sales data for the last four months:
		
		\vspace{4mm}
		\begin{center}
			\begin{tabular}{|c|c|}
				\hline
				\textbf{Month (\( t \))} & \textbf{Sales \(\$ Y_t \)} \\
				\hline
				\( t-3 \) & 120 \\
				\( t-2 \) & 130 \\
				\( t-1 \) & 140 \\
				\( t \)   & 150 \\
				\hline
			\end{tabular}
		\end{center}
		\vspace{4mm}
		
		\begin{enumerate}[label=(\alph*)]
			\item Calculate the 3-month Trailing Simple Moving Average (SMA) at time $t-1$ and \( t \).
%			\[
%			SMA_t = \frac{Y_{t-2} + Y_{t-1} + Y_t}{3}
%			\]
			\item What is the estimated trend value using the SMA method?

			\item Now, suppose we apply a Weighted Moving Average (WMA) where more recent months are given higher weights:
			\begin{itemize}
				\item Weight for \( t \) (most recent month) = 3
				\item Weight for \( t-1 \) (second most recent) = 2
				\item Weight for \( t-2 \) (third most recent) = 1
			\end{itemize}
			 Compute the 3-month WMA at time \( t-1 \) and \( t \) using the formula:
%			\[
%			WMA_t = \frac{w_{t} Y_t + w_{t-1} Y_{t-1} + w_{t-2} Y_{t-2}}{w_t + w_{t-1} + w_{t-2}}
%			\]
			
			\item How does the WMA estimate compare to the SMA estimate?
			\item Why does the Weighted Moving Average react more quickly to recent changes compared to the Simple Moving Average?
			\item Prof. Vince suggested us to use the following weight:
			\begin{itemize}
				\item Weight for \( t \) (most recent month) = 1
				\item Weight for \( t-1 \) (second most recent) = 1
				\item Weight for \( t-2 \) (third most recent) = 1
			\end{itemize}
			Is it different from the SMA? Does it do the job of ``more recent months are given higher weights"?
		\end{enumerate}
		
		\clearpage
		\textbf{Sol'n}
		\begin{enumerate}[label=(\alph*)]
			\item 
			The SMA at time \(t-1\) is given by:
			\[
			\text{SMA}_{t-1} = \frac{Y_{t-3} + Y_{t-2} + Y_{t-1}}{3} = \frac{120 + 130 + 140}{3} = \frac{390}{3} = 130.
			\]
			The SMA at time \(t\) is given by:
			\[
			\text{SMA}_t = \frac{Y_{t-2} + Y_{t-1} + Y_t}{3} = \frac{130 + 140 + 150}{3} = \frac{420}{3} = 140.
			\]
			
			\item 
			The trend estimate is taken as the SMA value, so the estimated trend at time $t-1$ is 130 and \(t\) is 140.

			\item 
			The WMA is computed by:
			\[
			\text{WMA}_{t-1} = \frac{w_0 Y_{t-1} + w_{-1} Y_{t-2} + w_{-2} Y_{t-3}}{w_0 + w_{-1} + w_{-2}}
			= \frac{3(140) + 2(130) + 1(120)}{3 + 2 + 1} = \frac{800}{6} \approx 133.33
			\]
			\[
			\text{WMA}_t = \frac{w_0 Y_t + w_{-1} Y_{t-1} + w_{-2} Y_{t-2}}{w_0 + w_{-1} + w_{-2}}
			= \frac{3(150) + 2(140) + 1(130)}{3 + 2 + 1} = \frac{860}{6} \approx 143.33
			\]
			
			\item 
			The SMA yields 140, whereas the WMA is approximately 143.33. The higher weight on the most recent month (with a sales figure of 150) causes the WMA to be higher.
			
			\item 
			Because the WMA assigns a larger weight to the most recent observation, it reacts more quickly to changes in recent data. In contrast, the SMA gives equal weight to all observations, smoothing out recent fluctuations.
			
			\item 
			With weights \(w_t = w_{t-1} = w_{t-2} = 1\), the WMA becomes:
			\[
			\text{WMA}_t = \frac{1\cdot Y_t + 1\cdot Y_{t-1} + 1\cdot Y_{t-2}}{1+1+1} = \frac{Y_t + Y_{t-1} + Y_{t-2}}{3},
			\]
			which is exactly the same as the SMA. Therefore, using equal weights does not emphasize recent observations.
			
		\end{enumerate}
		
		\clearpage
		
		
		
%		\newpage 
%		\centering{This page intentionally left blank.}
%		\ % The empty page
%		

%		 \vfill
%		 \centering{-- The End -- }
		
		
		
		%		\begin{solution}
			%			
			%		\end{solution}	
%		\clearpage 
%		
%		\question Was wäre, wenn es keine Luft gäbe?
%		\begin{parts} 
%			\part[5] Was würde mit Luftballons geschehen? 
%			\bonuspart[6] Wie könnten Fluggesellschaften damit umgehen?
%		\end{parts}
%		
%		\question [100] Wird es morgen schneien?
%		\begin{checkboxes}
%			\CorrectChoice Ja
%			\choice Nein
%			\choice Vielleicht
%		\end{checkboxes}
%		
%		\question Ein Name der folgenden Reihe passt nicht zu den anderen. Welcher?
%		\begin{oneparchoices}
%			\choice Donald
%			\choice Dagobert
%			\choice Daisy
%			\choice Micky
%			\CorrectChoice Balu
%		\end{oneparchoices}
		
	\end{questions}
	
%	\begin{center}
%		\combinedgradetable[h][questions]
%	\end{center}
	
\end{document}